% Generated draft: requires mathematical and source review.
\documentclass[11pt]{article}
\usepackage[margin=25mm]{geometry}
\usepackage{fontspec}
\setmainfont{Latin Modern Roman}
\usepackage{amsmath,amssymb,mathtools,unicode-math}
\setmathfont{Latin Modern Math}
\usepackage{enumitem,needspace,xurl,hyperref}
\setlength{\parindent}{0pt}
\setlength{\parskip}{5pt}
\setlength{\emergencystretch}{3em}
\newcommand{\N}{\mathbb{N}}
\newcommand{\Z}{\mathbb{Z}}
\newcommand{\Q}{\mathbb{Q}}
\newcommand{\R}{\mathbb{R}}
\newcommand{\C}{\mathbb{C}}
\newcommand{\F}{\mathbb{F}}
\newcommand{\A}{\mathbb{A}}
\newcommand{\PP}{\mathbb{P}}
\newcommand{\E}{\mathbb{E}}
\newcommand{\eps}{\varepsilon}
\newcommand{\dd}{\,\mathrm{d}}
\begin{document}
\section*{UnsolvedMath ID: 30002879}
\textbf{Source problem code:} catalogue entry
\subsection{Definitions and mathematical statement}
Write \(C_2=\Z/2\Z\). For a commutative unital base ring \(K\) of characteristic \(2\), put \(A=K[C_2\times C_2]\) and let \(M\) have underlying \(K\)-module \(K[C_2]\). To define the displayed triangular algebra, one must specify an \((A,K)\)-bimodule structure on \(M\), compatible with its \(K\)-module structure. With such a choice, define
\[
B=\begin{pmatrix}A&M\\0&K\end{pmatrix},\qquad
\begin{pmatrix}a&m\\0&\lambda\end{pmatrix}
\begin{pmatrix}a'&m'\\0&\lambda'\end{pmatrix}
=\begin{pmatrix}aa'&am'+m\lambda'\\0&\lambda\lambda'\end{pmatrix},
\qquad e=\begin{pmatrix}0&0\\0&1\end{pmatrix}.
\]
The source specifies only \(\operatorname{char}K=2\); it does not explicitly state that \(K\) is a field or give the left \(A\)-action on \(M\). The commutative-base convention above explains the notation and requires confirmation for the intended problem. In particular, no action induced by a particular homomorphism \(C_2\times C_2\to C_2\) is selected here.

The source assumes that \(e\) is stratifying: multiplication induces an isomorphism
\[
Be\otimes_{eBe}eB\longrightarrow BeB,
\]
and \(\operatorname{Tor}^{eBe}_i(Be,eB)=0\) for every \(i>0\). Here \(BeB\) is the two-sided ideal generated by \(e\).

Let \(\operatorname{HH}^n(B)=H^n(C^\bullet(B,B))\), where \(C^n(B,B)=\operatorname{Hom}_K(B^{\otimes_K n},B)\) carries the Hochschild differential. The standard Gerstenhaber bracket has degree \(-1\):
\[
[-,-]\colon\operatorname{HH}^p(B)\times\operatorname{HH}^q(B)\longrightarrow\operatorname{HH}^{p+q-1}(B).
\]
Thus its graded Lie interpretation uses shifted degree \(p-1\) on \(\operatorname{HH}^p(B)\).

\textbf{Requested statement, subject to the unresolved specifications.} For the intended base \(K\) and intended bimodule \(M\), determine the graded Lie structure of \(\operatorname{HH}^*(B)\) with its Gerstenhaber bracket, for example by a complete description of its graded components and brackets. Decide whether the intended long exact sequence associated with the short exact sequence of \(B\)-bimodules
\[
0\longrightarrow BeB\longrightarrow B\longrightarrow B/BeB\longrightarrow0
\]
determines that Lie structure. The source does not display the long exact sequence or specify its construction, terms, maps, or bracket compatibility; these must be identified before the second request is precise. The Gerstenhaber interpretation is the standard interpretation of the requested Lie structure, but the source does not explicitly specify a grading convention or request an additional restricted Lie operation. This draft records the problem and its unresolved qualifications without asserting a computation.
\subsection{Short English statement}
The problem asks for the Hochschild-cohomology Lie bracket of a triangular algebra in characteristic \(2\), and whether a long exact sequence determines that bracket. The base assumptions, the action defining the off-diagonal bimodule, and the intended long exact sequence need clarification. The supplied open status and verification labels are curation metadata, not independent mathematical verification.
\subsection{Sources}
\begin{itemize}
\item Local registry: \texttt{data/unsolvedmath/problems.json}, with \texttt{catalog/problems.json} as the additional-record registry; supplied complete record for permanent UnsolvedMath ID \(30002879\), exact source problem code \texttt{catalogue entry}, title ``Hochschild-Cohomology Lie Structure of a Stratified Algebra.'' Attribution is to the Oberwolfach Reports Open Problems collection; no individual proposer is supplied. The proposed year is recorded as \(2015\).
\item Supplied report citation: \emph{Cohomology of Finite Groups: Interactions and Applications} (2016), Oberwolfach Reports, pp.\ 1317--1384, DOI \texttt{10.4171/owr/2015/24}. \url{https://doi.org/10.4171/owr/2015/24}. The background describes these reports as documenting mathematical workshops and their open-problem discussions. The citation gives 2016 while its report identifier and proposed-year metadata refer to 2015; the bibliographic dating has not been independently checked.
\item Supplied literature source: Briggs--Rubio y Degrassi, \emph{Stable invariance of restricted Lie structure}. \url{https://arxiv.org/abs/2006.13871}. No more specific theorem or passage supporting this example is supplied.
\item Supplied literature source: Oberwolfach Report 2015/24. \url{https://doi.org/10.4171/owr/2015/24}.
\item Literature assessment dated exactly \texttt{2026-08-21}: the supplied curation labels the problem open and reports that the cited exact sequence does not by itself supply the requested Lie algebra computation, and that no complete computation of this example was found. This is supplied dated literature triage derived from the OWR curation, not an independently verified assessment; its status and cited claims require independent verification.
\item The supplied statement status is \texttt{corrected\_verified}, explained in the record as preservation of mathematical content while repairing notation, typography, or omitted definitions. This label is recorded as supplied metadata and does not establish mathematical correctness of this draft.
\end{itemize}
\end{document}
